%% 公文样例：通知
%% 构建：bash scripts/build_pdf.sh examples/notice.tex
\documentclass[UTF8, fontset=none, twoside]{ctexart}
\usepackage{gongwen}

\begin{document}

\gworg{华中科技大学教务处文件}
\gwdocno{教字〔2026〕15 号}
\gwtitle{关于开展数值方法课程资源建设的通知}

\gwto{各教学单位：}

\gwbody{为加强数值计算方法课程建设，提升学生运用计算机求解数学问题的能力，
经研究决定开展课程资源建设工作。现将有关事项通知如下。}

\gwhone{一、建设目标}

\gwhtwo{（一）总体要求}

\gwbody{以培养学生数值计算能力为核心，建设一套内容准确、体例统一、
可直接运行的中文教学资源。资源应覆盖浮点误差、方程求根、线性方程组、
插值、最小二乘、数值积分、常微分方程与偏微分方程等基本内容。}

\gwhtwo{（二）具体目标}

\gwhthree{1.\ 讲义建设}

\gwbody{编写中文讲义一部，正文不少于四十万字，配图不少于一百五十幅。
全部示例代码应可在 Jupyter 环境中直接运行，并给出可复核的数值结果。}

\gwhthree{2.\ 习题配套}

\gwbody{每章配习题六至十道，其中二至三道给出提示。习题应覆盖基本概念、
算法实现与误差分析三个层次。}

\gwhone{二、主要任务}

\gwhtwo{（一）内容编写}

\gwbody{按章分解任务，每章指定负责人一名，负责讲义撰写、代码实现与数值验证。
各章内容须经交叉复核，重点核查数值结论的正确性。}

\gwhtwo{（二）质量把关}

\gwbody{所有数值实验须与高精度结果或标准库对照，并在讲义中如实说明偏差。
凡不能复现的结论，一律不得写入讲义。}

\gwhone{三、进度安排}

\gwhtwo{（一）第一阶段}

\gwbody{二〇二六年十一月底前，完成讲义初稿并通过内部评议。}

\gwhtwo{（二）第二阶段}

\gwbody{二〇二六年十二月底前，完成修订、排版与印制，形成可分发版本。}

\gwhone{四、工作要求}

\gwbody{各教学单位要高度重视，统筹安排人员，保证编写时间。编写过程中
应注重体例统一，避免前后重复与表述不一致。遇有技术问题，及时与教务处联系。}

\gwhone{五、其他事项}

\gwhtwo{（一）经费保障}

\gwbody{课程建设经费从本科教学工程专项中列支，主要用于资料购置、
人员劳务与印制费用，具体额度另行核定。}

\gwhtwo{（二）联系方式}

\gwbody{联系人：教务处教学建设科。联系电话：000-00000000。}

\gwbody{本通知自印发之日起执行。}

% 附件与落款必须同页，否则会出现只有两行字的孤页
\begin{samepage}
\gwattach{附件：1.\ 课程资源建设任务分工表}
\gwattachcont{2.\ 讲义编写体例说明}
\gwsign{华中科技大学教务处}{二〇二六年十月二十六日}
\end{samepage}

\end{document}
